% GENERATED FILE. Edit canonical lessons, then run npm run generate:books.
% canonical-source-hash: fnv1a64:791b0922ff3255e1
% canonical-lessons: FR-C336-desaccord, FR-C336-objecter, FR-C336-objection, FR-C336-retorquer, FR-C336-contredire

\chapter{Pas d'accord}
\label{ch:fr-pas-d-accord}
\hlchaptermodality{336}

\begin{chapteropening}
\textbf{By the end of this chapter:} \emph{I can disagree politely: say we disagree, raise an objection, retort and contradict what someone said.}

Complete the last lesson of chapter 336: I can disagree politely: say we disagree, raise an objection, retort and contradict what someone said.
\end{chapteropening}

\section[le désaccord]{le désaccord — the disagreement}

\label{lesson:FR-C336-desaccord}



\begin{quote}
Before the new one: say the French for the line of argument, then the French for to persuade.
\end{quote}

\subsection*{You'll want to know: le désaccord}
\textbf{le désaccord} — \textquotedblleft{}the disagreement\textquotedblright{}.

\begin{quote}
Nous sommes en désaccord sur presque tout.
\end{quote}

\begin{quote}
\emph{We disagree about almost everything.}
\end{quote}

\emph{Être en désaccord} — to disagree. The opposite is \emph{être d'accord}.

\begin{cousinweb}
\emph{désaccord} is built from \emph{dés-} \textquotedblleft{}not\textquotedblright{} and \emph{accord} \textquotedblleft{}agreement\textquotedblright{}, which comes from the Old French verb \emph{acorder} (EtymDB). That verb is usually traced to a Latin formation on \emph{cor, cordis} \textquotedblleft{}heart\textquotedblright{}.
\end{cousinweb}

\begin{grammarlens}[title={What you've built}]
One more word for telling, explaining and describing at B1.
\end{grammarlens}

\subsection*{Your turn}
\begin{itemize}
\raggedright
  \item \emph{Say it:} \emph{le désaccord}
  \item \emph{Say it:} \emph{le désaccord}, once more
  \item \emph{Say it:} a sentence of your own with \emph{le désaccord}
  \item \emph{Recall:} say the French for the line of argument, then the French for to persuade, then say \emph{le désaccord} again
\end{itemize}

\begin{tcolorbox}[breakable,colback=teal!4,colframe=teal!35!black,title={Before you move on}]
What does le désaccord mean? (\textquotedblleft{}The disagreement\textquotedblright{}.) Say it once more.
\end{tcolorbox}
\section[objecter]{objecter — to object, to say in reply}

\label{lesson:FR-C336-objecter}



\begin{quote}
Before the new one: say the French for to persuade, then the French for the disagreement.
\end{quote}

\subsection*{You'll want to know: objecter}
\textbf{objecter} — \textquotedblleft{}to object, to say in reply\textquotedblright{}.

\begin{quote}
Elle a objecté que c'était trop cher.
\end{quote}

\begin{quote}
\emph{She objected that it was too expensive.}
\end{quote}

\emph{Objecter que…} — to object that… Formal; in speech many people say \emph{je ne suis pas d'accord}.

\begin{cousinweb}
\emph{objecter} is a learned borrowing from Latin \emph{obiectāre} \textquotedblleft{}to throw against, to set before\textquotedblright{} (Lewis \& Short), formed on \emph{obicere} \textquotedblleft{}to throw before\textquotedblright{}, from \emph{ob} and \emph{iacere} \textquotedblleft{}to throw\textquotedblright{}.
\end{cousinweb}

\begin{grammarlens}[title={What you've built}]
One more word for telling, explaining and describing at B1.
\end{grammarlens}

\subsection*{Your turn}
\begin{itemize}
\raggedright
  \item \emph{Say it:} \emph{objecter}
  \item \emph{Say it:} \emph{objecter}, once more
  \item \emph{Say it:} a sentence of your own with \emph{objecter}
  \item \emph{Recall:} say the French for to persuade, then the French for the disagreement, then say \emph{objecter} again
\end{itemize}

\begin{tcolorbox}[breakable,colback=teal!4,colframe=teal!35!black,title={Before you move on}]
What does objecter mean? (\textquotedblleft{}To object, to say in reply\textquotedblright{}.) Say it once more.
\end{tcolorbox}
\section[l'objection]{l'objection — the objection}

\label{lesson:FR-C336-objection}



\begin{quote}
Before the new one: say the French for the disagreement, then the French for to object.
\end{quote}

\subsection*{You'll want to know: l'objection}
\textbf{l'objection} — \textquotedblleft{}the objection\textquotedblright{}.

\begin{quote}
Pas d'objection ? Alors nous commençons.
\end{quote}

\begin{quote}
\emph{No objection? Then we start.}
\end{quote}

\emph{Faire une objection} — to raise an objection. Feminine.

\begin{cousinweb}
\emph{objection} is a learned borrowing from Late Latin \emph{obiectiō} \textquotedblleft{}a throwing against, an objection\textquotedblright{}, from \emph{obicere} \textquotedblleft{}to throw before\textquotedblright{} (Lewis \& Short), built on \emph{iacere} \textquotedblleft{}to throw\textquotedblright{}.
\end{cousinweb}

\begin{grammarlens}[title={What you've built}]
One more word for telling, explaining and describing at B1.
\end{grammarlens}

\subsection*{Your turn}
\begin{itemize}
\raggedright
  \item \emph{Say it:} \emph{l'objection}
  \item \emph{Say it:} \emph{l'objection}, once more
  \item \emph{Say it:} a sentence of your own with \emph{l'objection}
  \item \emph{Recall:} say the French for the disagreement, then the French for to object, then say \emph{l'objection} again
\end{itemize}

\begin{tcolorbox}[breakable,colback=teal!4,colframe=teal!35!black,title={Before you move on}]
What does l'objection mean? (\textquotedblleft{}The objection\textquotedblright{}.) Say it once more.
\end{tcolorbox}
\section[rétorquer]{rétorquer — to retort}

\label{lesson:FR-C336-retorquer}



\begin{quote}
Before the new one: say the French for to object, then the French for the objection.
\end{quote}

\subsection*{You'll want to know: rétorquer}
\textbf{rétorquer} — \textquotedblleft{}to retort\textquotedblright{}.

\begin{quote}
Il a rétorqué que ce n'était pas vrai.
\end{quote}

\begin{quote}
\emph{He retorted that it was not true.}
\end{quote}

Literally to twist back: you turn the words back on the speaker.

\begin{cousinweb}
\emph{rétorquer} is a learned borrowing from Latin \emph{retorquēre} \textquotedblleft{}to twist or bend back\textquotedblright{} (EtymDB; Lewis \& Short), built on \emph{torquēre} \textquotedblleft{}to twist\textquotedblright{}.
\end{cousinweb}

\begin{grammarlens}[title={What you've built}]
One more word for telling, explaining and describing at B1.
\end{grammarlens}

\subsection*{Your turn}
\begin{itemize}
\raggedright
  \item \emph{Say it:} \emph{rétorquer}
  \item \emph{Say it:} \emph{rétorquer}, once more
  \item \emph{Say it:} a sentence of your own with \emph{rétorquer}
  \item \emph{Recall:} say the French for to object, then the French for the objection, then say \emph{rétorquer} again
\end{itemize}

\begin{tcolorbox}[breakable,colback=teal!4,colframe=teal!35!black,title={Before you move on}]
What does rétorquer mean? (\textquotedblleft{}To retort\textquotedblright{}.) Say it once more.
\end{tcolorbox}
\section[contredire]{contredire — to contradict}

\label{lesson:FR-C336-contredire}



\begin{quote}
Before the new one: say the French for the objection, then the French for to retort.
\end{quote}

\subsection*{You'll want to know: contredire}
\textbf{contredire} — \textquotedblleft{}to contradict\textquotedblright{}.

\begin{quote}
Je n'aime pas quand tu me contredis devant tout le monde.
\end{quote}

\begin{quote}
\emph{I do not like it when you contradict me in front of everyone.}
\end{quote}

Like \emph{dire}, but \emph{vous contredisez}. \emph{Se contredire} — to contradict yourself.

\begin{cousinweb}
\emph{contredire} comes through Old French from Latin \emph{contrādīcere} \textquotedblleft{}to speak against\textquotedblright{} (EtymDB), built from \emph{contrā} \textquotedblleft{}against\textquotedblright{} and \emph{dīcere} \textquotedblleft{}to say\textquotedblright{} (Lewis \& Short).
\end{cousinweb}

\begin{grammarlens}[title={What you've built}]
One more word for telling, explaining and describing at B1.
\end{grammarlens}

\subsection*{Your turn}
\begin{itemize}
\raggedright
  \item \emph{Say it:} \emph{contredire}
  \item \emph{Say it:} \emph{contredire}, once more
  \item \emph{Say it:} a sentence of your own with \emph{contredire}
  \item \emph{Recall:} say the French for the objection, then the French for to retort, then say \emph{contredire} again
\end{itemize}

\begin{tcolorbox}[breakable,colback=teal!4,colframe=teal!35!black,title={Before you move on}]
What does contredire mean? (\textquotedblleft{}To contradict\textquotedblright{}.) Say it once more.
\end{tcolorbox}
